\originalchapter{8.04 作业7：双点势与谐振子}
来源为Spring 2016 Assignment 7, Barton Zwiebach. 原件第1、3、4、6题完整, 本篇中文转述与\textbf{AI独立解答}. 第2题Griffiths 2.47、p.87, 第5题Griffiths 2.42、p.86仅有题号, 不补造. 第2题附MIT建议将上一作业模型作数值检验, 已在作业6数值解中实施, 不宣称取得外书题面.
\originalsection{第1题：双点势束缚与分裂}
\begin{exercise}
$V=-g[\delta(x-a)+\delta(x+a)]$,$g>0$, $\lambda=mag/\hbar^2$. 求偶、奇束缚态超越方程, 以$\hbar^2/(ma^2)$为能量单位作图; 大距离时求分裂.
\end{exercise}
\begin{solution}
$E=-\hbar^2\kappa^2/(2m)$, 中间偶态为$C\cosh\kappa x$, 奇态为$C\sinh\kappa x$, 两端为衰减指数. $x=a$连续与导数跳跃$\psi'(a+)-\psi'(a-)=-2mg\psi(a)/\hbar^2$给$s=\lambda(1+\e^{-2s})$（偶）、$s=\lambda(1-\e^{-2s})$（奇）,$s=\kappa a$. 偶态任意$\lambda>0$都有一根, 奇态严格$\lambda>1/2$才有非零根, 其零根不能计为束缚. 能量为$-s^2/2$乘题给单位, 曲线见下图.

令$\kappa_0=mg/\hbar^2$, 大$\lambda=\kappa_0a$时$\kappa_\pm\simeq\kappa_0(1\pm\e^{-2\kappa_0a})$. 因而$E_o-E_e\simeq2mg^2\e^{-2mga/\hbar^2}/\hbar^2$. 偶态更低, 分裂随两中心距离$2a$指数降低; 指数里不能再把$a$误当完整距离.
\begin{center}\begin{tikzpicture}
\begin{axis}[width=.82\textwidth,height=4.2cm,xlabel={$\lambda$},ylabel={$E/[\hbar^2/(ma^2)]$},grid=major,legend style={font=\small},legend pos=south west]
\addplot+[mark=none] coordinates{(0.01,-0.000192379796) (0.0224583333,-0.000927226455) (0.0349166667,-0.00214797659) (0.047375,-0.00379934747) (0.0598333333,-0.00583625237) (0.0722916667,-0.0082214598) (0.08475,-0.0109238905) (0.0972083333,-0.0139173543) (0.109666667,-0.0171795958) (0.122125,-0.0206915637) (0.134583333,-0.0244368417) (0.147041667,-0.0284012015) (0.1595,-0.0325722468) (0.171958333,-0.0369391266) (0.184416667,-0.0414923032) (0.196875,-0.0462233622) (0.209333333,-0.0511248551) (0.221791667,-0.0561901695) (0.23425,-0.0614134198) (0.246708333,-0.0667893558) (0.259166667,-0.0723132844) (0.271625,-0.077981004) (0.284083333,-0.0837887479) (0.296541667,-0.0897331354) (0.309,-0.0958111301) (0.321458333,-0.102020003) (0.333916667,-0.108357302) (0.346375,-0.114820823) (0.358833333,-0.121408585) (0.371291667,-0.12811881) (0.38375,-0.134949906) (0.396208333,-0.141900446) (0.408666667,-0.148969154) (0.421125,-0.156154897) (0.433583333,-0.163456665) (0.446041667,-0.170873567) (0.4585,-0.178404818) (0.470958333,-0.186049733) (0.483416667,-0.193807716) (0.495875,-0.201678256) (0.508333333,-0.209660919) (0.520791667,-0.217755343) (0.53325,-0.225961232) (0.545708333,-0.234278352) (0.558166667,-0.242706528) (0.570625,-0.251245634) (0.583083333,-0.2598956) (0.595541667,-0.268656396) (0.608,-0.27752804) (0.620458333,-0.286510588) (0.632916667,-0.295604134) (0.645375,-0.304808809) (0.657833333,-0.314124774) (0.670291667,-0.323552224) (0.68275,-0.333091382) (0.695208333,-0.342742497) (0.707666667,-0.352505844) (0.720125,-0.362381723) (0.732583333,-0.372370455) (0.745041667,-0.382472382) (0.7575,-0.392687867) (0.769958333,-0.403017291) (0.782416667,-0.41346105) (0.794875,-0.42401956) (0.807333333,-0.434693248) (0.819791667,-0.445482559) (0.83225,-0.456387948) (0.844708333,-0.467409885) (0.857166667,-0.478548849) (0.869625,-0.489805332) (0.882083333,-0.501179834) (0.894541667,-0.512672865) (0.907,-0.524284943) (0.919458333,-0.536016595) (0.931916667,-0.547868355) (0.944375,-0.559840762) (0.956833333,-0.571934363) (0.969291667,-0.584149708) (0.98175,-0.596487356) (0.994208333,-0.608947865) (1.00666667,-0.621531802) (1.019125,-0.634239735) (1.03158333,-0.647072234) (1.04404167,-0.660029873) (1.0565,-0.673113229) (1.06895833,-0.686322878) (1.08141667,-0.6996594) (1.093875,-0.713123374) (1.10633333,-0.726715381) (1.11879167,-0.740436) (1.13125,-0.754285812) (1.14370833,-0.768265397) (1.15616667,-0.782375334) (1.168625,-0.7966162) (1.18108333,-0.810988571) (1.19354167,-0.825493023) (1.206,-0.840130127) (1.21845833,-0.854900453) (1.23091667,-0.86980457) (1.243375,-0.884843042) (1.25583333,-0.90001643) (1.26829167,-0.915325293) (1.28075,-0.930770187) (1.29320833,-0.946351661) (1.30566667,-0.962070262) (1.318125,-0.977926535) (1.33058333,-0.993921017) (1.34304167,-1.01005424) (1.3555,-1.02632674) (1.36795833,-1.04273903) (1.38041667,-1.05929164) (1.392875,-1.07598508) (1.40533333,-1.09281985) (1.41779167,-1.10979647) (1.43025,-1.12691542) (1.44270833,-1.14417721) (1.45516667,-1.1615823) (1.467625,-1.17913119) (1.48008333,-1.19682435) (1.49254167,-1.21466224) (1.505,-1.23264532) (1.51745833,-1.25077406) (1.52991667,-1.26904888) (1.542375,-1.28747024) (1.55483333,-1.30603857) (1.56729167,-1.3247543) (1.57975,-1.34361785) (1.59220833,-1.36262962) (1.60466667,-1.38179003) (1.617125,-1.40109948) (1.62958333,-1.42055835) (1.64204167,-1.44016704) (1.6545,-1.45992592) (1.66695833,-1.47983536) (1.67941667,-1.49989573) (1.691875,-1.52010738) (1.70433333,-1.54047067) (1.71679167,-1.56098593) (1.72925,-1.5816535) (1.74170833,-1.60247371) (1.75416667,-1.62344687) (1.766625,-1.64457331) (1.77908333,-1.66585333) (1.79154167,-1.68728723) (1.804,-1.7088753) (1.81645833,-1.73061782) (1.82891667,-1.75251508) (1.841375,-1.77456735) (1.85383333,-1.79677489) (1.86629167,-1.81913795) (1.87875,-1.8416568) (1.89120833,-1.86433167) (1.90366667,-1.8871628) (1.916125,-1.91015042) (1.92858333,-1.93329476) (1.94104167,-1.95659603) (1.9535,-1.98005444) (1.96595833,-2.00367021) (1.97841667,-2.02744352) (1.990875,-2.05137457) (2.00333333,-2.07546354) (2.01579167,-2.09971062) (2.02825,-2.12411598) (2.04070833,-2.14867978) (2.05316667,-2.1734022) (2.065625,-2.19828337) (2.07808333,-2.22332347) (2.09054167,-2.24852262) (2.103,-2.27388097) (2.11545833,-2.29939865) (2.12791667,-2.3250758) (2.140375,-2.35091253) (2.15283333,-2.37690896) (2.16529167,-2.40306521) (2.17775,-2.42938138) (2.19020833,-2.45585758) (2.20266667,-2.4824939) (2.215125,-2.50929044) (2.22758333,-2.53624729) (2.24004167,-2.56336453) (2.2525,-2.59064224) (2.26495833,-2.61808051) (2.27741667,-2.64567938) (2.289875,-2.67343895) (2.30233333,-2.70135927) (2.31479167,-2.72944039) (2.32725,-2.75768237) (2.33970833,-2.78608528) (2.35216667,-2.81464914) (2.364625,-2.84337401) (2.37708333,-2.87225993) (2.38954167,-2.90130694) (2.402,-2.93051506) (2.41445833,-2.95988433) (2.42691667,-2.98941478) (2.439375,-3.01910643) (2.45183333,-3.0489593) (2.46429167,-3.07897342) (2.47675,-3.10914878) (2.48920833,-3.13948542) (2.50166667,-3.16998334) (2.514125,-3.20064254) (2.52658333,-3.23146303) (2.53904167,-3.26244482) (2.5515,-3.2935879) (2.56395833,-3.32489228) (2.57641667,-3.35635794) (2.588875,-3.38798488) (2.60133333,-3.41977309) (2.61379167,-3.45172256) (2.62625,-3.48383329) (2.63870833,-3.51610524) (2.65116667,-3.54853842) (2.663625,-3.58113279) (2.67608333,-3.61388835) (2.68854167,-3.64680506) (2.701,-3.67988291) (2.71345833,-3.71312188) (2.72591667,-3.74652193) (2.738375,-3.78008303) (2.75083333,-3.81380517) (2.76329167,-3.84768831) (2.77575,-3.88173242) (2.78820833,-3.91593746) (2.80066667,-3.95030341) (2.813125,-3.98483023) (2.82558333,-4.01951789) (2.83804167,-4.05436635) (2.8505,-4.08937556) (2.86295833,-4.1245455) (2.87541667,-4.15987613) (2.887875,-4.1953674) (2.90033333,-4.23101928) (2.91279167,-4.26683172) (2.92525,-4.30280468) (2.93770833,-4.33893813) (2.95016667,-4.37523201) (2.962625,-4.41168628) (2.97508333,-4.44830091) (2.98754167,-4.48507585) (3,-4.52201104)};
\addlegendentry{偶态}
\addplot+[mark=none] coordinates{(0.508333333,-0.000137366812) (0.520791667,-0.000841420667) (0.53325,-0.00211823943) (0.545708333,-0.00394184536) (0.558166667,-0.0062882168) (0.570625,-0.00913510428) (0.583083333,-0.0124618673) (0.595541667,-0.0162493296) (0.608,-0.0204796492) (0.620458333,-0.0251362038) (0.632916667,-0.0302034866) (0.645375,-0.0356670135) (0.657833333,-0.0415132399) (0.670291667,-0.0477294851) (0.68275,-0.0543038643) (0.695208333,-0.0612252269) (0.707666667,-0.068483101) (0.720125,-0.0760676424) (0.732583333,-0.0839695887) (0.745041667,-0.0921802171) (0.7575,-0.100691306) (0.769958333,-0.1094951) (0.782416667,-0.118584278) (0.794875,-0.127951923) (0.807333333,-0.137591497) (0.819791667,-0.147496813) (0.83225,-0.157662015) (0.844708333,-0.168081555) (0.857166667,-0.178750175) (0.869625,-0.189662889) (0.882083333,-0.200814965) (0.894541667,-0.212201909) (0.907,-0.223819454) (0.919458333,-0.235663545) (0.931916667,-0.247730325) (0.944375,-0.260016126) (0.956833333,-0.272517457) (0.969291667,-0.285230994) (0.98175,-0.298153573) (0.994208333,-0.311282178) (1.00666667,-0.324613935) (1.019125,-0.338146103) (1.03158333,-0.351876069) (1.04404167,-0.365801341) (1.0565,-0.379919539) (1.06895833,-0.394228392) (1.08141667,-0.408725732) (1.093875,-0.423409489) (1.10633333,-0.438277685) (1.11879167,-0.45332843) (1.13125,-0.46855992) (1.14370833,-0.483970428) (1.15616667,-0.499558305) (1.168625,-0.515321975) (1.18108333,-0.531259929) (1.19354167,-0.547370725) (1.206,-0.563652985) (1.21845833,-0.580105391) (1.23091667,-0.59672668) (1.243375,-0.613515646) (1.25583333,-0.630471135) (1.26829167,-0.647592043) (1.28075,-0.664877313) (1.29320833,-0.682325935) (1.30566667,-0.699936942) (1.318125,-0.717709411) (1.33058333,-0.735642456) (1.34304167,-0.753735231) (1.3555,-0.771986928) (1.36795833,-0.790396774) (1.38041667,-0.808964028) (1.392875,-0.827687983) (1.40533333,-0.846567964) (1.41779167,-0.865603324) (1.43025,-0.884793448) (1.44270833,-0.904137744) (1.45516667,-0.92363565) (1.467625,-0.943286628) (1.48008333,-0.963090166) (1.49254167,-0.983045773) (1.505,-1.00315298) (1.51745833,-1.02341135) (1.52991667,-1.04382045) (1.542375,-1.06437988) (1.55483333,-1.08508925) (1.56729167,-1.1059482) (1.57975,-1.12695638) (1.59220833,-1.14811345) (1.60466667,-1.1694191) (1.617125,-1.19087304) (1.62958333,-1.21247498) (1.64204167,-1.23422465) (1.6545,-1.25612179) (1.66695833,-1.27816617) (1.67941667,-1.30035755) (1.691875,-1.32269573) (1.70433333,-1.34518049) (1.71679167,-1.36781165) (1.72925,-1.39058903) (1.74170833,-1.41351245) (1.75416667,-1.43658176) (1.766625,-1.45979681) (1.77908333,-1.48315747) (1.79154167,-1.50666359) (1.804,-1.53031507) (1.81645833,-1.55411179) (1.82891667,-1.57805364) (1.841375,-1.60214054) (1.85383333,-1.62637238) (1.86629167,-1.65074911) (1.87875,-1.67527063) (1.89120833,-1.69993688) (1.90366667,-1.72474782) (1.916125,-1.74970337) (1.92858333,-1.7748035) (1.94104167,-1.80004816) (1.9535,-1.82543733) (1.96595833,-1.85097096) (1.97841667,-1.87664904) (1.990875,-1.90247155) (2.00333333,-1.92843847) (2.01579167,-1.95454979) (2.02825,-1.98080551) (2.04070833,-2.00720562) (2.05316667,-2.03375014) (2.065625,-2.06043906) (2.07808333,-2.0872724) (2.09054167,-2.11425018) (2.103,-2.1413724) (2.11545833,-2.16863911) (2.12791667,-2.19605031) (2.140375,-2.22360603) (2.15283333,-2.25130632) (2.16529167,-2.2791512) (2.17775,-2.30714071) (2.19020833,-2.33527488) (2.20266667,-2.36355377) (2.215125,-2.39197741) (2.22758333,-2.42054586) (2.24004167,-2.44925915) (2.2525,-2.47811735) (2.26495833,-2.5071205) (2.27741667,-2.53626866) (2.289875,-2.56556188) (2.30233333,-2.59500023) (2.31479167,-2.62458376) (2.32725,-2.65431253) (2.33970833,-2.68418662) (2.35216667,-2.71420607) (2.364625,-2.74437096) (2.37708333,-2.77468136) (2.38954167,-2.80513733) (2.402,-2.83573893) (2.41445833,-2.86648625) (2.42691667,-2.89737936) (2.439375,-2.92841831) (2.45183333,-2.95960319) (2.46429167,-2.99093408) (2.47675,-3.02241103) (2.48920833,-3.05403414) (2.50166667,-3.08580347) (2.514125,-3.1177191) (2.52658333,-3.14978111) (2.53904167,-3.18198957) (2.5515,-3.21434457) (2.56395833,-3.24684618) (2.57641667,-3.27949447) (2.588875,-3.31228953) (2.60133333,-3.34523144) (2.61379167,-3.37832027) (2.62625,-3.41155611) (2.63870833,-3.44493903) (2.65116667,-3.47846911) (2.663625,-3.51214644) (2.67608333,-3.54597109) (2.68854167,-3.57994314) (2.701,-3.61406267) (2.71345833,-3.64832976) (2.72591667,-3.68274449) (2.738375,-3.71730695) (2.75083333,-3.7520172) (2.76329167,-3.78687533) (2.77575,-3.82188142) (2.78820833,-3.85703555) (2.80066667,-3.89233779) (2.813125,-3.92778822) (2.82558333,-3.96338693) (2.83804167,-3.99913398) (2.8505,-4.03502947) (2.86295833,-4.07107345) (2.87541667,-4.10726602) (2.887875,-4.14360725) (2.90033333,-4.18009721) (2.91279167,-4.21673598) (2.92525,-4.25352363) (2.93770833,-4.29046025) (2.95016667,-4.3275459) (2.962625,-4.36478066) (2.97508333,-4.40216461) (2.98754167,-4.43969781) (3,-4.47738034)};
\addlegendentry{奇态}
\end{axis}\end{tikzpicture}\end{center}

\end{solution}
\originalsection{第3题：单侧经典禁区概率}
\begin{exercise}
谐振子的$n=0,1,2$态, 求位置大于正的经典振幅$A_n$的概率.
\end{exercise}
\begin{solution}
令$d=\sqrt{\hbar/(m\omega)}$, $\xi=x/d$, $A_n/d=\sqrt{2n+1}$. 分布$\rho_n(\xi)=H_n^2(\xi)\e^{-\xi^2}/(\sqrt\pi 2^n n!)$, 从正转折点积分到无穷. 令$z_n=\sqrt{2n+1}$, $\operatorname{erfc}z=2\int_z^\infty\e^{-u^2}\dd u/\sqrt\pi$, 得
\[
P_0=\tfrac12\operatorname{erfc}(1),\quad
P_1=\tfrac12\operatorname{erfc}(\sqrt3)+\frac{\sqrt3\e^{-3}}{\sqrt\pi},\quad
P_2=\tfrac12\operatorname{erfc}(\sqrt5)+\frac{\e^{-5}11\sqrt5}{2\sqrt\pi}.
\]
分部积分$\int_z^\infty u^{2r}\e^{-u^2}\dd u=z^{2r-1}\e^{-z^2}/2+(2r-1)\int_z^\infty u^{2r-2}\e^{-u^2}\dd u/2$可逐项核验. 数值分别0.0786496035、0.0558051125、0.0475347174. 这是题问的单侧; 若问$|x|>A_n$才加倍. 禁区概率非零不是经典转折点失效, 而是量子态没有确定经典轨道.
\end{solution}
\originalsection{第4题：四次矩与Hermite生成函数}
\begin{exercise}
(a)用升降算符求$\langle x^4\rangle_n$. (b)分别求第$n$个能量本征态的$\Delta x$与$\Delta p$, 再求乘积$\Delta x\Delta p$. (c)从$\e^{-s^2+2s\xi}=\sum H_n(\xi)s^n/n!$证明最高次项与Hermite微分方程.
\end{exercise}
\begin{solution}
$x=d(a+a^\dagger)/\sqrt2$. 为避免不对易展开, 先作用一次$x^2$:
$x^2\ket n=(d^2/2)[\sqrt{n(n-1)}\ket{n-2}+(2n+1)\ket n+\sqrt{(n+1)(n+2)}\ket{n+2}]$. 取范数平方即$\langle x^4\rangle=3(d^2/2)^2(2n^2+2n+1)$. 两均值0,$\langle x^2\rangle=d^2(n+1/2)$,$\langle p^2\rangle=\hbar^2(n+1/2)/d^2$, 故$\Delta x=d\sqrt{n+1/2}$, $\Delta p=\hbar\sqrt{n+1/2}/d$, $\Delta x\Delta p=\hbar(n+1/2)$, 仅基态饱和基本界.

生成函数$G$中$s^n$最高$\xi$次项来自$\e^{2s\xi}$的$n$项, 乘$n!$得$(2\xi)^n$. 直接求导有$G_{\xi\xi}=4s^2G$, $G_\xi=2sG$, $G_s=(-2s+2\xi)G$, 故$G_{\xi\xi}-2\xi G_\xi+2sG_s=0$. 比较$s^n/n!$系数得$H_n''-2\xi H_n'+2nH_n=0$. 生成函数提供系数关系的证明, 不是只列已知多项式.
\end{solution}
\originalsection{第6题：两个能态的摆动}
\begin{exercise}
初态$(\varphi_0+\varphi_1)/\sqrt2$, 求时间波函数、密度、位置和动量均值. 一般谐振子叠加态的密度周期是什么?
\end{exercise}
\begin{solution}
$\psi(t)=[\varphi_0\e^{-\ii\omega t/2}+\varphi_1\e^{-3\ii\omega t/2}]/\sqrt2$. 实的本征函数给$\rho=(\varphi_0^2+\varphi_1^2+2\varphi_0\varphi_1\cos\omega t)/2$. 因$\bra0x\ket1=d/\sqrt2$, 有$\langle x\rangle=d\cos\omega t/\sqrt2$, $\langle p\rangle=-\hbar\sin\omega t/(\sqrt2d)$, 与$\dot{\langle x\rangle}=\langle p\rangle/m$相合. 位置均值摆动振幅$d/\sqrt2$, 角频率$\omega$; 动量幅为$\hbar/(\sqrt2d)$. 一般态各$n$在$T=2\pi/\omega$演化都乘$-1$, 所以密度周期为$T$或其约数; 单能態密度定常, 不能断言所有态的最小周期都恰为$T$.
\end{solution}
